% sec-02: the action register.  Table 17, Figure 11 (plantuml-type activity
% diagram with a fork, a join, and a connector).
\section{The Action Register}

Table 17 lists the day's seven actions in the order they are done. The answer to
the inquiry goes first, because every other letter refers to it. The Pitch waits
on the organization record, and the letter to NIH waits on the Pitch, because it
reports the Pitch as submitted.

\begin{apptable}
\begin{tabularx}{\textwidth}{>{\raggedright\arraybackslash}p{0.5cm} >{\raggedright\arraybackslash}p{6.4cm} >{\raggedright\arraybackslash}p{2.6cm} >{\raggedright\arraybackslash}X}
\headrow \thead{No} & \thead{Action} & \thead{File} & \thead{Waits on}\\
\midrule
1 & Answer the NSF inquiry & {\footnotesize\ttfamily email-01} & Nothing \\
2 & Ask the pilot's gate of its investigator & {\footnotesize\ttfamily email-02} & Row 1 sent \\
3 & Confirm the Research.gov record & {\footnotesize\ttfamily form-02} & An active SAM.gov record \\
4 & Submit the Project Pitch & {\footnotesize\ttfamily form-01} & Row 3 \\
5 & Ask about NAIRR Pilot resources & {\footnotesize\ttfamily email-03} & Nothing \\
6 & Disclose the Pitch to NIH & {\footnotesize\ttfamily email-04} & Row 4 submitted \\
7 & Open the match ledger & {\footnotesize\ttfamily capital-04} & Nothing \\
\end{tabularx}
\end{apptable}
\vspace{-0.60cm}
\tabcap{Table 17. The day's seven actions in the order they are done, each with the file that\\
carries it and what it waits on; the letter to NIH follows the Pitch it discloses.}

\subsection{From the Pitch to a proposal}

Figure 11 draws the route from the answer to a proposal as an activity diagram.
Two branches run in parallel and must join before the Pitch is submitted: the
registration branch, which confirms the company's federal entity record and its
Research.gov organization \cite{researchgov2026}, and the drafting branch, which
writes the four Pitch fields and measures each against its character limit. The
diagram breaks at a connector and resumes at the right. NSF then invites a full
proposal or declines to \cite{nsf26510}. A declined Pitch is revised and
submitted again, and the solicitation allows two Pitches a year.

\begin{appfloat}[!tb]
\begin{appfig}[plantuml-type / activity diagram with a fork, a join, and a connector]
% Left block: the answer, a fork into two branches, the join, and the Pitch,
% ending at connector A.  Right block: connector A, the NIH disclosure, the
% one decision, and the stop, with a note on what an award would open.
\node[umlinit] (st) at (3.50,0.30) {};
\node[umlbox,text width=46mm] (a1) at (3.50,-0.35) {Answer the NSF inquiry, letter 1};
\node[umlbar,minimum width=70mm] (fk) at (3.50,-1.00) {};
\node[umlbox,text width=40mm] (b1) at (1.25,-1.65) {Confirm SAM.gov and the UEI};
\node[umlbox,text width=40mm] (b2) at (1.25,-2.60) {Research.gov organization};
\node[umlbox,text width=40mm] (b3) at (1.25,-3.55) {Assign the PI and AOR roles};
\node[umlbox,text width=40mm] (c1) at (5.75,-1.65) {Draft the four Pitch fields};
\node[umlbox,text width=40mm] (c2) at (5.75,-2.60) {Check each field's limit};
\node[umlbox,text width=40mm] (c3) at (5.75,-3.55) {Confirm no field proposes a trial};
\node[umlbar,minimum width=70mm] (jn) at (3.50,-4.20) {};
\node[umlkey,text width=46mm] (sb) at (3.50,-4.85) {Submit the Project Pitch};
\node[circle,draw=fundink,line width=0.6pt,fill=fundwhite,minimum size=4.2mm,
  inner sep=0pt,font=\tiny\sffamily\bfseries] (ca) at (3.50,-5.55) {A};
\draw[umlarrow] (st) -- (a1);
\draw[umlarrow] (a1) -- (fk);
\draw[umlarrow] (fk.south -| b1) -- (b1);
\draw[umlarrow] (fk.south -| c1) -- (c1);
\draw[umlarrow] (b1) -- (b2);
\draw[umlarrow] (b2) -- (b3);
\draw[umlarrow] (c1) -- (c2);
\draw[umlarrow] (c2) -- (c3);
\draw[umlarrow] (b3) -- (jn.north -| b3);
\draw[umlarrow] (c3) -- (jn.north -| c3);
\draw[umlarrow] (jn) -- (sb);
\draw[umlarrow] (sb) -- (ca);
\node[circle,draw=fundink,line width=0.6pt,fill=fundwhite,minimum size=4.2mm,
  inner sep=0pt,font=\tiny\sffamily\bfseries] (cb) at (10.60,0.30) {A};
\node[umlbox,text width=46mm] (d1) at (10.60,-0.35) {Disclose the Pitch to NIH, letter 4};
\node[diamond,aspect=2.4,draw=fundink,line width=0.6pt,fill=fundwhite,inner sep=1pt,
  text width=26mm,font=\tiny\sffamily,align=center] (dc) at (10.60,-1.55) {Invited to submit a proposal?};
\node[umlbox,text width=40mm] (p1) at (10.60,-2.95) {Proposal by the next deadline};
\node[umlbox,text width=40mm] (p2) at (10.60,-3.90) {Merit review on NSF's criteria};
\node[umlbox,text width=24mm] (r1) at (14.35,-2.95) {Revise and pitch again};
\node[diamond,draw=fundink,line width=0.6pt,fill=fundwhite,minimum size=3.6mm,inner sep=0pt] (mg) at (10.60,-4.70) {};
\node[umlfinal] (en) at (10.60,-5.45) {};
\fill[fundink] (en) circle (0.9mm);
\node[umlnote,text width=27mm] (nt) at (14.10,-5.45) {Only a Phase II award opens the pilot and Strategic Breakthrough};
\draw[umlarrow] (cb) -- (d1);
\draw[umlarrow] (d1) -- (dc);
\draw[umlarrow] (dc.south) -- node[umlguard,right,pos=0.35] {[yes]} (p1.north);
\draw[umlarrow] (dc.east) -| node[umlguard,above,pos=0.25] {[no]} (r1.north);
\draw[umlarrow] (p1) -- (p2);
\draw[umlarrow] (p2) -- (mg);
\draw[umlarrow] (r1.south) |- (mg.east);
\draw[umlarrow] (mg) -- (en);
\draw[umldash,-] (nt.west) -- (en.east);
\ptitle{-0.75}{0.95}{Two branches join before the Pitch}
\end{appfig}
\vspace{-0.60cm}
\figcaption{Figure 11. The route from the answer to a proposal, with registration and drafting\\
run in parallel and joined before the Pitch, then one decision that NSF alone makes.}
\end{appfloat}

\subsection{What each letter asks}

Letter 1 thanks NSF's SBIR program for the inquiry, states the gate, and asks
whether the inquiry contemplates a route the company has missed. It copies the
Policy Office, which owns the rules the gate is written in, and the Research.gov
help desk, which would resolve any registration fault \cite{nsfseedcontact,nsfpappg}.
Letter 2 asks the pilot's principal investigator at the University of Central
Florida to confirm the gate in writing, and asks for nothing else
\cite{nctiucf2026}. Letter 3 asks the NAIRR Pilot whether a small business can
request shared resources for verification research on synthetic data only
\cite{nsfnairr2026,nairrpilot2026}. Letter 4 tells the NIH SBIR program office
and the National Cancer Institute's SBIR team that an NSF Pitch has been
submitted on the verification work, so that any later NIH application discloses
it and neither agency pays for the same work twice
\cite{nihsbirseed,ncisbir2026,sbapolicydirective}.

\subsection{A note on the date}

October 1 is the first day of the federal fiscal year. If appropriations have
lapsed, NSF staff may not answer for a time, and a portal may accept a
submission it cannot yet process. The Pitch can still be submitted. Nothing else
is followed up until three business days after the agency reopens.

\actionbox{The one approval this day asks for}{Approve answering the NSF inquiry,
testing the pilot's published eligibility gate in writing with its principal
investigator, and submitting the Project Pitch that opens the only route to it.
Approve disclosing the Pitch to NIH on the day it is submitted.}
